\chapter*{参考文献}\addcontentsline{toc}{chapter}{参考文献}\markboth{参考文献}{参考文献}
\section*{缩写标注原始文献}

\BibEntry{AP1: Stumpf, Kilian, edited by Paul Rule, Claudia von Collani. The Acta Pekinensia, or, Historical Records of the Maillard de Tournon Legation: First Transcribed Edition and English Annotated Translation. Rome : Institutum Historicum Societatis Iesu; Macau: Macau Ricci Institute, 2015.}

\BibEntry{AP2: Stumpf, Kilian, edited by Rule, Paul and Collani, Claudia von. The Acta Pekinensia, or, Historical Records of the Maillard de Tournon Legation. Volume II: September 1706 –December 1707, Leiden: Brill, 2019.}

\BibEntry{AP3: Stumpf, Kilian, edited by Rule, Paul and Collani, Claudia von. The Acta Pekinensia or Historical Records of the Maillard de Tournon Legation, Volume III: January 1708 –February 1709, Leiden: Brill, 2024.}

\section*{其他原始文献（按文献名音序排列）}

\BibEntry{An Account of Tibet: The Travels of Ippolito Desideri of Pistoia, S. J., 1712–1727, Desideri, Ippolito, London: Routledge, 1932.}

\BibEntry{Letters of a Peking Jesuit: The Correspondence of Ferdinand Verbiest, SJ (1623-1688) , Leuven: Ferdinand Verbiest Institute, University of Leuven, 2017.}

\BibEntry{Travels from St. Petersburg, in Russia, to Diverse Parts of Asia, Bell, John, Glasgow: printed for the author by R. and A. Foulis, 1763.}

\BibEntry{《本朝题驳公案》，《清代判牍案例汇编》（甲编），册1，北京：社会科学文献出版社，2019。}

\BibEntry{《澄怀园文存》，张廷玉，《近代中国史料丛刊》第52辑，台北：文海出版社，1970。}

\BibEntry{《鞑靼旅行记》，《清代西人见闻录》，北京：中国人民大学出版社，1985。}

\BibEntry{《大清律例》，《景印文渊阁四库全书》，台北：台湾商务印书馆，1986。}

\BibEntry{《大义觉迷录》，内务府武英殿本，哈佛燕京珍本网络版，编号：T4661.823243。}

\BibEntry{《戴铎口供》，《〈文献丛编〉全编》，第3册，北京：北京图书馆出版社，2008。}

\BibEntry{《憺园文集》，徐乾学，冠山堂刻本（共36卷）。}

\BibEntry{《定藏纪程》，《川藏游踪汇编》，成都：四川民族出版社，1985。}

\BibEntry{《读书堂西征随笔》，北京：国立故宫博物院，1928。}

\BibEntry{《方苞集》，上海：上海古籍出版社，1983。}

\BibEntry{《抚远大将军允禵奏稿》，吴丰培编，北京：全国图书馆文献微缩复制中心，1991。}

\BibEntry{《高澹人行述》，清沈氏啸园蓝格钞本，\url{https://rbook.ncl.edu.tw}。}

\BibEntry{《关于江宁织造曹家档案史料》，北京：中华书局，1975。}

\BibEntry{《皇清钦授征仕郎日讲官起居注翰林院检讨祖考竹垞府君行述》，《广清碑传集》，钱仲联，苏州：苏州大学出版社，1999。}

\BibEntry{《已未词科录》，秦瀛，哈佛大学藏1807年刻本（编号：4668.81679），\url{https://catalog.hathitrust.org/Record/011635144}。}

\BibEntry{《郡主觉罗氏墓志铭》，《纳兰家族墓碑》，《景印文渊阁四库全书》，台北：台湾商务印书馆，1986。}

\BibEntry{《康熙朝汉文朱批奏折汇编》，北京：中国档案出版社，1984。}

\BibEntry{《康熙朝满文朱批奏折全译》，北京：中国社会科学出版社，1996。}

\BibEntry{《康熙三十二年俄罗斯商人义迭思〈聘盟日记〉》，《历史档案》，2004，第4期。}

\BibEntry{《揆方墓志铭》，《纳兰家族墓碑》，《景印文渊阁四库全书》，台北：台湾商务印书馆，1986。}

\BibEntry{《李文贞公年谱》，《近代中国史料丛刊》第63辑，台北：文海出版社，1970。}

\BibEntry{《笠翁一家言全集》，李渔，芥子园刻本，哈佛燕京珍本网络版，编号：T5499.4433。}

\BibEntry{《穆经远口供》，《〈文献丛编〉全编》，第2册，北京：北京图书出版公司，2008。}

\BibEntry{《年羹尧满汉奏折译编》，天津：天津古籍出版社，1995。}

\BibEntry{《平定三逆方略》，《景印文渊阁四库全书》，台北：台湾商务印书馆，1986。}

\BibEntry{《平定准噶尔方略》，《景印文渊阁四库全书》，台北：台湾商务印书馆，1986。}

\BibEntry{《平湖县志》，乾隆五十五年刻本，哈佛燕京珍本网络版，编号：990080818430203941。}

\BibEntry{《曝书亭全集》，朱彝尊，长春：吉林文史出版社，2009。}

\BibEntry{《清初西洋传教士满文档案译本》，郑州：大象出版社，2015。}

\BibEntry{《清代起居注册·康熙朝》，北京、台北：中华书局、联经出版事业公司联合出版，2009。}

\BibEntry{《清史稿》，北京：中华书局，1976。}

\BibEntry{《榕村语录榕村语录续集》，李光地，北京：商务印书馆，2019。}

\BibEntry{《圣祖仁皇帝亲征平定朔漠方略》，《景印文渊阁四库全书》，台北：台湾商务印书馆，1986。}

\BibEntry{《圣祖仁皇帝圣训》，《景印文渊阁四库全书》，台北：台湾商务印书馆，1986。}

\BibEntry{《圣祖仁皇帝实录》，《大清历朝实录》，台北：华联出版社，1964。}

\BibEntry{《圣祖仁皇帝庭训格言》，《景印文渊阁四库全书》，台北：台湾商务印书馆，1986。}

\BibEntry{《世宗皇帝上谕内阁》，《景印文渊阁四库全书》，台北：台湾商务印书馆，1986。}

\BibEntry{《世宗宪皇帝上谕八旗》，《景印文渊阁四库全书》，台北：台湾商务印书馆，1986。}

\BibEntry{《世宗宪皇帝实录》，《大清历朝实录》，台北：华联出版社，1964。}

\BibEntry{《世宗宪皇帝朱批谕旨》，《景印文渊阁四库全书》，台北：台湾商务印书馆，1986。}

\BibEntry{《世祖章皇帝实录》，《大清历朝实录》，台北：华联出版社，1964。}

\BibEntry{《唐明律合编》，薛允升编，北京：法律出版社，1999。}

\BibEntry{《万寿庆典初集》，《景印文渊阁四库全书》，台北：台湾商务印书馆，1986。}

\BibEntry{《〈文献丛编〉全编》，北京：北京图书馆出版社，2008。}

\BibEntry{《啸亭杂录》，昭梿，北京：中华书局，1980。}

\BibEntry{《言鲭》，吕种玉，康熙四十四年刻本，早稻田大学图书馆（Waseda UniversityLibrary）藏，扫描版本编号：04003470018。}

\BibEntry{《耶稣会士中国书简集中国回忆录》，郑州：大象出版社，2001。}

\BibEntry{《雍正朝汉文朱批奏折汇编》，南京：江苏古籍出版社，1989。}

\BibEntry{《雍正朝满文朱批奏折全译》，合肥：黄山书社，1998。}

\BibEntry{《雍正朝起居注册》，中国第一历史档案馆编，北京：中华书局，1993。}

\BibEntry{《永宪录》，萧奭，北京：中华书局，1959。}

\BibEntry{《与顾梁汾书》，《通志堂集》，纳兰性德，上海：上海古籍出版社，1979。}

\BibEntry{《元以来西藏地方与中央政府关系档案史料汇编》，中国藏学研究中心等编，北京：中国藏学出版社，1994。}

\BibEntry{《允禩允禟案》，《〈文献丛编〉全编》，第2册，北京：北京图书馆出版社，2008。}

\BibEntry{《允禩允禟案续》，《〈文献丛编〉全编》，第3册，北京：北京图书馆出版社，2008。}

\BibEntry{《詹事府少詹事查公墓志铭》，《海宁查氏家族文化研究》，杭州：浙江大学出版社，2006。}

\BibEntry{《子平真诠评注》，沈孝瞻撰、徐乐吾评注、赵嘉宁注译，北京：中医古籍出版社，2012。}

\BibEntry{《奏请逆犯汪景祺等枯骨由菜市口竹签移下掩埋折》，雍正十三年十月二十九日，故宫008633号，台北故宫博物院清代档案检索系统。}

\BibEntry{《奏请赵显忠补靖远协副将折》，雍正六年五月二十一日，故宫003252号，台北故宫博物院清代档案检索系统。研究论著}

\BibEntry{Bailey, Michael D., "The Disenchantment of Magic: Spells, Charms, and Superstition in Early European Witchcraft Literature." The American Historical Review, 111, no. 2, 2006.}

\BibEntry{Bartlett, Beatrice, Monarchs and Ministers: The Grand Council in Mid-Ch'ing China, 1723-1820 , Berkeley: University of California Press, 1990.}

\BibEntry{Brian C. Jayne, Fred E. Inbau, and Joseph P. Buckley, Criminal Interrogation and Confessions , Burlington: Jones \& Bartlett Learning, 2013.}

\BibEntry{Chan, Albert, "A European Document on the Fall of the Ming Dynasty (1644–1649),"Monumenta Serica, vol. 35 (1981–1983).}

\BibEntry{Collani, Claudia von, "Kilian Stumpf and his Acta Pekinensia' : Life, Context and Purpose of a Manuscript," in Macau Ricci Institute eds., Acta Pekinensia: Western Historical Sources for the Kangxi Reign: International Symposium , Macau: Macau Ricci Institute, 2013.}

\BibEntry{Collingwood, R.G., The Idea of History, Oxford: Oxford University Press, 1956.}

\BibEntry{Croce, Benedetto, Aesthetic: As Science of Expression and General Linguistic, London: Macmillan, 1909.}

\BibEntry{Croce, Benedetto, History as The Story of Liberty, London: George Allen and Unwin Limited, 1941.}

\BibEntry{Croce, Benedetto, "History Brought Under the General Concept of Art" Lecture, Accademia Pontaniana from University of Naples, Naples, March 5, 1893, translated and edited by Brian P. Copenhaver, Rebecca Copenhaver, From Kant to Croce: Modern Philosophy in Italy, 1800–1950 ,Toronto: University of Toronto Press, 2012.}

\BibEntry{Curtis, Emily Byrne, "Aspects of a Multi-Faceted Process: The Circulation of Enamel Wares between the Vatican and Kangxi's Court (1700–1722)" , Extreme-Orient Extreme-Occident, vol.43, no.1, 2019.}

\BibEntry{Dai, Yingcong, The Sichuan Frontier and Tibet: Imperial Strategy in the Early Qing, Seattle: University of Washington Press, 2011.}

\BibEntry{Domanska, Ewa, Encounters: Philosophy of History after Postmodernism, Charlottesville: University of Virginia Press, 1998.}

\BibEntry{Elman, Benjamin A. On Their Own Terms: Science in China, 1550–1900, Cambridge, MA: Harvard University Press, 2005.}

\BibEntry{Feng, Jia, "The Emperor's Coffer: The Qing Imperial Fiscal Separation Between Privy Purse and State Treasury (1644-1912)." PhD diss., University of California, Los Angeles, 2017.}

\BibEntry{Fishman, Joshua, Language Loyalty in the United States: The Maintenance and Perpetuation of Non-English Mother Tongues by American Ethnic and Religious Groups, The Hague: Mouton \& Co., 1966.}

\BibEntry{G. Županov, Ines, and Pierre Antoine Fabre, eds., The Rites Controversies in the Early Modern World, Leiden: Brill, 2018.}

\BibEntry{Hellman, Lisa, "Enslaved in Dzungaria: What an Eighteenth-Century Crocheting Instructor Can Teach Us about Overland Globalization.' Journal of Global History, vol. 17, no. 3, 2022.}

\BibEntry{Herodotus, Rawlinson, George trans, History, London: John Murray, 1880.}

\BibEntry{Huang, Pei, Reorienting the Manchus: A Study of Sinicization, 1583–1795, Ithaca: Cornell University Press, 2011.}

\BibEntry{Huang, Pei, Autocracy at Work: A Study of the Yung-cheng Period: 1723-1735, Bloomington: Indiana University Press, 1976.}

\BibEntry{Jin, Guoping, "The Mission to Congratulate the Kangxi Emperor on his Seventeenth Birthday: Father Joao Mourao and His Plan to Bring an Ambassador of the King of Portugal to Beijing," in Macau Ricci Institute eds., Acta Pekinensia: Western Historical Sources for the Kangxi Reign: International Symposium.}

\BibEntry{Kessler, Lawrence D., K'ang-hsi and the Consolidation of Ch'ing Rule, 1661–1684, Chicago: University of Chicago Press, 1976.}

\BibEntry{Minamiki, George, The Chinese Rites Controversy from Its Beginning to Modern Times, Chicago: Loyola University Press, 1985.}

\BibEntry{Novick, Peter, That Noble Dream: The "Objectivity Question" and the American Historical Profession, Cambridge: Cambridge University Press, 1988.}

\BibEntry{Pace, Edward A., et al., The Catholic Encyclopedia : An International Work of Reference on the Constitution, Doctrine, Discipline, and History of the Catholic Church, New York : Encyclopedia Press, 1913.}

\BibEntry{Petech, Luciano, "Notes on Tibetan History of the 18th Century", T’oung Pao , vol. 52, no. 4/5, 1966.}

\BibEntry{Petech, Luciano, China and Tibet in the Early 18th Century: History of Establishment of Chinese Protectorate in Tibet, Leiden: Brill, 1950.}

\BibEntry{Pih, Irene, and Le P. Gabriel de Magalhães, Un jésuite portugais en Chine au XVII e siècle, Paris: Fundação Calouste Gulbenkian-Centro Cultural Português, 1979.}

\BibEntry{Plato, The Republic, Cambridge: Cambridge University Press, 2000.}

\BibEntry{Poppe, Nicholas, "Renat's Kalmuck Maps", Imago Mundi 12 (1955):157-59.}

\BibEntry{Puente-Ballesteros, Beatriz, "Jesuit Medicine in the Kangxi Court (1662–1722): Imperial Networks and Patronage", East Asian Science, Technology, and Medicine, no.34, 2011.}

\BibEntry{Perdue, Peter, China Marches West: The Qing Conquest of Central Eurasia, Cambridge, MA: Harvard University Press, 2005.}

\BibEntry{Reed, Bradly W., Talons and Teeth: County Clerks and Runners in the Qing Dynasty, Stanford: Stanford University Press, 2000.}

\BibEntry{Rossabi, Morris, "The Tea and Horse Trade with Inner Asia During the Ming." Journal of Asian History, vol. 4, no. 2, 1970.}

\BibEntry{Rule, Paul, "The Acta Pekinensia: A Key Source for the History of Christianity in China”，古伟瀛编，《东西交流史的新局：以基督宗教为中心》，台北：台湾大学出版社，2005。}

\BibEntry{Shaughnessy, Edward, " 'New' Evidence on the Zhou Conquest", Early China, vol. 6, 1980.}

\BibEntry{Spence, Jonathan, Emperor of China, New York: Vintages Book, 1988.}

\BibEntry{Spence, Jonathan, Ts'ao Yin and the K'ang-hsi Emperor: Bondservant and Master, New Haven: Yale University Press, 1966.}

\BibEntry{Storozum, Michael, Lu, P., Wang, S. et al. "Geoarchaeological Evidence of the AD 1642 Yellow River Flood That Destroyed Kaifeng, a Former Capital of Dynastic China". Scientific Reports 10, 3765 (2020).}

\BibEntry{Sweet, Michael J. , and Zwilling, Leonard, "The First Capuchin Mission to Tibet: Fr. Domenico da Fano's Report of 1713." Buddhist-Christian Studies, vol. 42, 2022.}

\BibEntry{Swen, Litian, "Bypassing the Bureaucracy: Qing Rulership and the First Papal Legation to China", Late Imperial China, 2024, v.2.}

\BibEntry{Thiers, Consuelo, and Wehner, Leslie, "The Personality Traits of Populist Leaders and Their Foreign Policies: Hugo Chávez and Donald Trump" International Studies Quarterly, vol.66, Issue 1, 2022.}

\BibEntry{Tucker, K. H., Jr. , " A New Kind of Personal Politics: On the Populism of Trump and Sanders", Journal of Political Ideologies, 30(1), 2022.}

\BibEntry{Vannini,Fulgentius, The Bell of Lhasa , New Delhi: Messelas Devarsons, 1976.}

\BibEntry{Wakeman, Frederic, The Fall of Imperial China, New York: The Free Press, 1975.}

\BibEntry{White, Hayden, Metahistory: The Historical Imagination in Nineteenth- Century Europe, Johns Hopkins University Press, 2014.}

\BibEntry{White, Hayden, The Content of the Form: Narrative Discourse and Historical Representation, Baltimore Johns Hopkins University Press, 1987.}

\BibEntry{Witek, John W., Controversial Ideas in China and in Europe: a Biography of Jean-François Foucquet, S.J., (1665-1741) , Roma: Institutum Historicum S.I., 1982.}

\BibEntry{Wittgenstein, Ludwig, G. E. M. Anscombe trans., Philosophical Investigations, London: Basil Blackwell, 1958.}

\BibEntry{Wu, Marshall, “A-er-hsi-p'u and His Painting Collection" , in Li, Chu- tsing ed., Artists and Patrons: Some Social and Economic Aspects of Chinese Paintings, Seattle: University of Washington Press, 1989.}

\BibEntry{Wu, Silas H. L., Communication and Imperial Control in China: Evolution of the Palace Memorial system, Cambridge, MA: Harvard University Press, 1970.}

\BibEntry{Wu, Silas H. L., Passage to Power: K'ang-his and His Heir Apparent, 1661–1722, Cambridge, MA: Harvard University Press, 1979.}

\BibEntry{Zelin, Madeleine. "Yung-Cheng Reign," in The Cambridge History of China: The Ch'ing Empire to 1800. vol.9. Cambridge: Cambridge University Press, 2002.}

\BibEntry{白新良，《清史考辨》，北京：人民出版社，2006。}

\BibEntry{白新良，《清史考辨》，北京：中华书局，2021。}

\BibEntry{柏拉图著，郭斌和、张竹明译，《理想国》，北京：商务印书馆，1986。}

\BibEntry{蔡崇榜，《宋代修史制度研究》，台北：文津出版社，1993。}

\BibEntry{蔡元培，《〈石头记〉索隐》，上海：商务印书馆，1934。}

\BibEntry{蔡志纯，《论蒙藏封建主围绕真假六世达赖喇嘛展开的争权斗争》，《中国蒙古史学会论文选集（1983）》，呼和浩特：内蒙古人民出版社，1983。}

\BibEntry{常建华，《从西方文献看雍正帝继位的合法性》，《天津师范大学学报》，2017，第2期。}

\BibEntry{常江、李理，《清宫大内侍卫》，北京：故宫出版社，2013。}

\BibEntry{陈葆真，《康熙和乾隆二帝的南巡及其对江南名胜和园林的绘制与仿建》，《故宫学术季刊》，2015，第32卷第3期。}

\BibEntry{陈葆真，《康熙皇帝的生日礼物与相关问题之探讨》，《故宫学术季刊》，2012，第30卷第1期。}

\BibEntry{陈国栋，《康熙小臣养心殿总监造赵昌生平小考》，冯明珠编，《盛清社会与扬州研究》，台北：远流出版公司，2011。}

\BibEntry{陈捷先，《康熙皇帝与书法》，《故宫学术季刊》，1999，第17卷第1期。}

\BibEntry{陈捷先，《雍正初年清世宗与年羹尧之君臣关系》，《台大历史学报》，1978，第5期。}

\BibEntry{陈捷先，《雍正写真》，北京：商务印书馆，2011。}

\BibEntry{陈寅恪，《金明馆丛稿二编》，北京：生活·读书·新知三联书店，2001。}

\BibEntry{陈寅恪，《唐代政治史述论稿》，北京：生活·读书·新知三联书店，2001。}

\BibEntry{陈寅恪，《元白诗笺证稿》，北京：生活·读书·新知三联书店，2001。}

\BibEntry{陈垣，《钱名世轶事》，《陈垣学术论文集》第二集，北京：中华书局，1982。}

\BibEntry{陈祖武，《清初学术思辨录》，北京：中国社会科学出版社，1992。}

\BibEntry{邓锐龄，《1720年率军进入拉萨的清军将领——延信》，《邓锐龄藏族史论文译文集》，北京：中国藏学出版社，2004。}

\BibEntry{邓锐龄，《1720年清军进入西藏的经过》，《邓锐龄藏族史论文译文集》，北京：中国藏学出版社，2004。}

\BibEntry{邓锐龄，《拉萨布达拉山东侧康熙时汉文摩崖考释》，《邓锐龄藏族史论文译文集》，北京：中国藏学出版社，2004。}

\BibEntry{邓锐龄，《年羹尧在雍正朝初期治藏政策孕育过程中的作用》，《邓锐龄藏族史论文译文集》，北京：中国藏学出版社，2004。}

\BibEntry{定宜庄，《满汉文化交流史话》，北京：社会科学文献出版社，2011。}

\BibEntry{董少新，《浅谈全球史的史料问题：以明清鼎革史的西文原始史料为中心》，《首都师范大学学报》，2018，第5期。}

\BibEntry{董少新，《形神之间：早期西洋医学入华史稿》，上海：上海古籍出版社，2008。}

\BibEntry{杜家骥，《雍正帝继位前的封旗及相关问题考析》，《中国史研究》，1990，第4期。}

\BibEntry{杜家骥，《雍正之弟改名阿其那、塞思黑问题试析》，《满族研究》，1998，第2期。}

\BibEntry{方甦生，《清太祖实录纂修考》，《辅仁学志》，1938，第7卷。}

\BibEntry{方志远，《明代国家权力结构及运行机制》，桂林：广西师范大学出版社，2024。}

\BibEntry{房兆楹，杜联喆，《增校清朝进士题名碑录附引得》，台北：成文书局，1966。}

\BibEntry{冯尔康，《清史史料学》，北京：故宫出版社，2013。}

\BibEntry{冯尔康，《雍正传》，上海：上海三联书店，1999。}

\BibEntry{冯尔康，《雍正继位新探》，天津：天津人民出版社，2008。}

\BibEntry{高学姣，《〈大义觉迷录〉的文本流播与记忆再现》，《清史研究》，2022，第5期。}

\BibEntry{顾诚，《论钦定〈明史〉作〈奸臣传〉之微意——兼为严嵩等人重新定位》，《顾诚明清史文集》，北京：北京日报出版社，2023。}

\BibEntry{关雪玲，《康熙时期西洋医学在清宫中传播问题的再考察》，《故宫学术季刊》，2012，第1期。}

\BibEntry{郭福祥，《巴多明与清宫西洋学馆》，《紫禁城》，1997，第4期。}

\BibEntry{郭沫若，《中国古代社会研究》，北京：商务印书馆，2011。}

\BibEntry{侯天印、王春华编著，《痰证论》，北京：人民军医出版社，1989。}

\BibEntry{胡春丽，《国图藏孤本〈抱耒堂集〉作者汪耀麟生平钩沉》，《古籍整理研究学刊》，2024，第2期。}

\BibEntry{胡琦，《已未词科与清初“文”“学”之辨》，《北京大学学报（哲学社会科学版）》，2014，第9期。}

\BibEntry{胡祥雨，《清代通奸罪文及其演变》，《历史研究》，2025，第8期。}

\BibEntry{胡小伟，《纳兰词到〈红楼梦〉——满汉文化融合的“青萍之争”》，《明清论丛》，1999，第一辑，北京：紫禁城出版社。}

\BibEntry{黄爱平，《毛奇龄与明末清初的学术》，《清史研究》，1996，第4期。}

\BibEntry{黄天骥，《纳兰性德和他的词》，广州：广东人民出版社，1983。}

\BibEntry{简欣晨，《康熙皇帝的集字活动—从几件董其昌（款）〈古诗十九首〉谈起》，《故宫学术季刊》，2021，第38卷第4期。}

\BibEntry{金承艺，《清朝帝位之争史事考》，北京：中华书局，2010。}

\BibEntry{金恒源，《雍正称帝与其对手》，上海：上海人民出版社，2008。}

\BibEntry{金启琮，《清朝前期卫拉特蒙古和中原的互市》，《内蒙古大学学报》，1964，第2期。}

\BibEntry{孔令伟，《钦差喇嘛楚儿沁藏布兰木占巴、清代西藏地图测绘与世界地理知识之传播》，《“中央研究院”历史语言研究所集刊》，2021，第3期。}

\BibEntry{李保文，《孝庄太皇太后与康熙朝〈泥金写本藏文龙藏经〉》，《故宫文物月刊》，卷335，2011。}

\BibEntry{李宪庆、白新良，《康、雍之际继嗣制度的演变》，《社会科学辑刊》，1983，第2期。}

\BibEntry{励俊，《追寻江村秘藏：高士奇书画鉴藏与流散研究》，杭州：中国美术学院出版社，2021。}

\BibEntry{梁绍杰，《清康熙朝大学士王掞奏请建储始末》，《故宫博物院院刊》，2007，第1期。}

\BibEntry{铃木真，《关于镶白旗和硕雍亲王胤禛与藩邸旧人》，《满学研究》，第7辑，北京：民族出版社，2002。}

\BibEntry{刘恒，《中国书法史：清代卷》，南京：江苏教育出版社，1999。}

\BibEntry{刘金库，《南画北渡：清代书画鉴藏中心研究》，石家庄：河北教育出版社，2008。}

\BibEntry{刘世珣，《底野迦的再现：康熙年间宫廷西药德里鸦噶初探》，《清史研究》，2014，第3期。}

\BibEntry{刘文鹏，《清代驿传及其与疆域形成关系之研究》，北京：中国人民大学出版社，2004。}

\BibEntry{刘小萌，《关于清代北京的俄罗斯人八旗满洲俄罗斯佐领历史寻踪》，《清史论丛》，2007年号，北京：中国广播电视出版社，2006。}

\BibEntry{刘小萌，《乾隆帝恢复允禩允禟宗籍考》，《历史档案》，2019，第2期。}

\BibEntry{刘小萌，《清代北京旗人社会》，北京：中国社会科学出版社，2016。}

\BibEntry{刘亚刚，《清初古书画“北移”现象之辨》，《文艺研究》，2018，第9期。}

\BibEntry{柳陞祺，《十八世纪初清政府平定西藏准噶尔之乱始末》，《民族研究》，1998，第1期。}

\BibEntry{陆蓓容，《王鸿绪尺牍考释—管窥清代书画鉴藏史（上）》，《中国美术学院学报》，2019，第3期。}

\BibEntry{毛亦可，《康熙末年的“面谕”与第三次立储计划——以〈本朝题驳公案〉新见史料为中心》，《清史研究》，2023，第5期。}

\BibEntry{茅海建，《天朝的崩溃：鸦片战争再研究》，北京：生活·读书·新知三联书店，2014。}

\BibEntry{孟妹芳，《蔡斑与年羹尧案关系考略》，《清史研究》，2005，第1期。}

\BibEntry{孟森，《读〈清实录〉商榷》，《明清史论著集刊》，北京：中华书局，1959。}

\BibEntry{孟森，《己未词科录外录》，《明清史论著集刊》，北京：中华书局1959。}

\BibEntry{孟森，《康熙重修太祖实录跋》，《明清史论著集刊》，北京：中华书局，1959。}

\BibEntry{孟森，《明清史讲义》，北京：中华书局，1981。}

\BibEntry{孟森，《清世宗入承大统考实》，《明清史论著集刊》，北京：中华书局，1959。}

\BibEntry{祁美琴、强光美编译，《满文〈满洲实录〉译编》，北京：中国人民大学出版社，2015。}

\BibEntry{乔治忠，《清朝官方史学研究》，台北：文津出版社，1994。}

\BibEntry{若松宽，《策妄阿喇布坦的崛起》，《清代蒙古的历史与宗教》，哈尔滨：黑龙江教育出版社，1994。}

\BibEntry{沈原，《“阿其那”、“塞思黑”考释》，《清史研究》，1997，第1期。}

\BibEntry{石岩刚，《〈御制平定西藏碑〉校录》，《西北民族论丛》，第7辑，北京：中国社会科学出版社，2010。}

\BibEntry{史松，《论年羹尧之死》，《清史研究》，1991，第3期。}

\BibEntry{史松，《雍正研究》，沈阳：辽宁民族出版社，2009。}

\BibEntry{宋嗣喜，《策妄阿喇布坦与沙皇俄国博罗库尔干出使俄国、温科夫斯基出使准噶尔与达尔扎复又出使俄国评述》，《中国蒙古史学会论文选集（1983）》，呼和浩特：内蒙古人民出版社，1983。}

\BibEntry{孙立天，《康熙的红票：全球化中的清朝》，北京：商务印书馆，2024。}

\BibEntry{汤开建，《沉与浮：明清鼎革变局中的欧洲传教士利类思与安文思》（上、下），《北京行政学院学报》，2014，第4、5期。}

\BibEntry{王光越，《“雍亲王谕”考释》，《清史研究》，2007，第3期。}

\BibEntry{王记录，《清代史馆与清代政治》，北京：人民出版社，2009。}

\BibEntry{王俊中，《“满洲”与“文殊”的渊源及西藏政教思想中的领袖与佛菩萨》，《“中央研究院”近代史研究所集刊》，卷28，1997。}

\BibEntry{王冕森，《康熙朝第二次废太子事件新探以新见托合齐口供档案为中心》，《清史论丛》，第48辑，2025。}

\BibEntry{王冕森，《清代后妃杂识》，上海：上海社会科学院出版社，2022。}

\BibEntry{王明珂，《华夏边缘：历史记忆与族群认同》，台北：允晨文化，1997。}

\BibEntry{王明珂，《羌在汉藏之间：川西羌族的历史人类学研究》，北京：中华书局，2008。}

\BibEntry{王佩环，《从新发现的满文档案再释阿其那与塞思黑》，《故宫博物院院刊》，2000，第2期。}

\BibEntry{王书奴，《中国娼妓史》，上海：生活书店，1934。}

\BibEntry{王树林，《高士奇年谱》，杭州：浙江古籍出版社，2021。}

\BibEntry{王思治，《索额图其人兼论助皇太子“潜谋大事”》，《清史研究》，1992，第1期。}

\BibEntry{王锺翰，《〈清史稿〉与清史研究》，《王锺翰清史论集》，北京：中华书局，2004。}

\BibEntry{王锺翰，《清圣祖遗诏考辨》，《王锺翰清史论集》，北京：中华书局，2004。}

\BibEntry{王锺翰，《清世宗夺嫡考实》，此文最早刊载于1949年《燕京学报》，《王锺翰清史论集》，北京：中华书局，2004。}

\BibEntry{王锺翰，《三释阿其那与塞思黑》，《王锺翰清史论集》，北京：中华书局，2004。}

\BibEntry{王锺翰，《胤祯西征纪实》，《王锺翰清史论集》，北京：中华书局，2004。}

\BibEntry{乌云毕力格，《五色四藩：多语文本中的内亚民族史地研究》，上海：上海古籍出版社，2016。}

\BibEntry{吴秀良，《允禵更名与雍正继位问题再探讨》，《清史研究》，2013，第3期。}

\BibEntry{武佩圣，《王翚客京师期间之交往与绘画活动》，《清初四王画派研究：“四王绘画艺术国际研讨会”论文集》，上海：上海书画出版社，1993。}

\BibEntry{夏伯嘉，《从安文思与汤若望之间的恩怨论说明宫廷耶稣会士》，《北京行政学院学报》，2017，第3期。}

\BibEntry{谢贵安，《清实录研究》，上海：上海古籍出版社，2013。}

\BibEntry{许曾重，《清世宗胤禛继承皇位问题新探》，《清史论丛》，1983，第4辑。}

\BibEntry{薛吕，《清宫造办处玻璃厂、玻璃作名称及厂址考证》，《故宫博物院院刊》，2023，第4期。}

\BibEntry{杨晋龙，《王应麟现存论著在清前期经学专著的传播考征（I）——以〈四库全书·经部〉收录书籍为对象的探讨》，《新经学》，第4辑，上海：上海人民出版社，2019。}

\BibEntry{杨启樵，《雍正篡位说驳难》，上海：上海书店出版社，2012。}

\BibEntry{杨启樵，《雍正帝及其密折制度研究》，上海：上海古籍出版社，2003。}

\BibEntry{杨珍，《废太子允礽一家与咸安宫》，《紫禁城》，2010，第7期。}

\BibEntry{杨珍，《关于康熙朝储位之争及雍正继位的几个问题》，《清史论丛》，1985，第6辑。}

\BibEntry{杨珍，《关于满文档案与清史研究的几点认识》，《故宫学刊》，2012。}

\BibEntry{杨珍，《康熙皇帝一家》，第五版，北京：学苑出版社，2021。}

\BibEntry{杨珍，《满文密折所见诚亲王允祉与雍正帝胤禛》，《第一届中日学者中国古代史论坛文集》，北京：中国社会科学出版社，2010。}

\BibEntry{杨珍，《清朝皇位继承制度》，第三版，北京：学苑出版社，2021。}

\BibEntry{杨珍，《清前期宫廷政治释疑》，北京：中国社会科学出版社，2018。}

\BibEntry{杨珍，《权臣索额图的身世与家事》，《紫禁城》，1996，第4期。}

\BibEntry{杨珍，《苏麻喇姑与康熙帝》，《故宫博物院院刊》，1995，第1期。}

\BibEntry{杨珍，《索额图研究》，《清史论丛》，1996，第13辑。}

\BibEntry{杨珍，《允禵储君地位问题研究》，《清史论丛》，1992，第9辑。}

\BibEntry{杨珍，《允禵与小呼毕勒罕的交往》，《紫禁城》，1991，第1期。}

\BibEntry{姚念慈，《康熙盛世与帝王心术：评“自古得天下之正莫如我朝”》，北京：生活·读书·新知三联书店，2018。}

\BibEntry{叶高树，《康熙皇帝的焦虑—立储、开边和历史定位》，《台湾师大历史学报》，2015，卷54。}

\BibEntry{叶高树，《满洲亲贵与清初政治：都英额地方赫舍里家族的个案研究》，《台湾师大历史学报》，2010，卷43。}

\BibEntry{雍琦，《朱彝尊年谱》，复旦大学硕士论文，2007。}

\BibEntry{尤旭东，《遥领芝告向东藩—清朝出使朝鲜使臣研究》，《故宫博物院院刊》，2014年，第6期。}

\BibEntry{余子铨，《“名教罪人”案与常州钱氏家族》，香港大学博士论文，2020。}

\BibEntry{袁鳞，《从捐纳到荐举：三藩之乱中的制科事件与士人心态——兼论已未词科的诏试缘由》，《苏州大学学报（哲学社会科学版）》，2024，第4期。}

\BibEntry{张建，《再造强权—准噶尔珲台吉策妄阿喇布坦崛起史新探》，《“中央研究院”历史语言研究所集刊》，2015，第86本第1分。}

\BibEntry{张勉治著，刘文鹏、王珏译，《康熙皇帝首次南巡与文武价值观念之间的对立》，《清史研究》，2011，第1期。}

\BibEntry{张书才，《年羹尧生年考实》，《历史档案》，1989，第1期。}

\BibEntry{张先清，《帝国潜流：清代前期的天主教、底层秩序与生活世界》，北京：社会科学文献出版社，2021。}

\BibEntry{赵柄学，《策妄阿拉布坦和清朝关系史之研究—从〈清内阁蒙古堂档〉研究18世纪清朝和准噶尔蒙古关系》，《中国边疆民族研究》，2011。}

\BibEntry{赵刚，《康熙博学鸿词科与清初政治变迁》，《故宫博物院院刊》，1993，第1期。}

\BibEntry{郑天挺，《清代包衣制度与宦官》，《探微集》，北京：中华书局，1980。}

\BibEntry{郑小悠，《年羹尧之死》，太原：山西人民出版社，2018。}

\BibEntry{郑小悠、橘玄雅、夏天，《九王夺嫡》，太原：山西人民出版社，2021。}

\BibEntry{钟国豪，《清初文臣高士奇（1645—1703）绘画鉴藏品味研究》，“中央大学”硕士论文，2022年。}

\BibEntry{朱金甫，《论康熙时期的南书房》，《故宫博物院院刊》，1990，第2期。}

\BibEntry{朱金甫，张书才主编，《清代典章制度辞典》，北京：中国人民大学出版社，2011。}

\BibEntry{朱维铮编，《周予同经学史论著选集》，上海：上海人民出版社，1983。}

\BibEntry{朱文锋主编，《中医诊断学》，上海：上海科学技术出版社，1995。}

\BibEntry{庄吉发，《从故宫旧档看清实录的窜改》，《清史拾遗》，台北：学生书局，1992。}

\BibEntry{庄吉发，《清代史料论述（一）》，台北：文史哲出版社，1979。}

\BibEntry{庄吉发，《清世宗拘禁十四阿哥胤始末》，《大陆杂志》，1974，第49卷，第2期。}